\documentclass[11pt,a4paper]{article}
\usepackage[top=1.65cm,bottom=1.7cm,left=1.8cm,right=1.8cm,headheight=14pt,footskip=20pt]{geometry}
\usepackage[T1]{fontenc}
\usepackage[utf8]{inputenc}
\usepackage[spanish]{babel}
\usepackage{lmodern,microtype,amsmath,xcolor,enumitem,titlesec,fancyhdr,hyperref}
\definecolor{ink}{HTML}{173F35}
\definecolor{muted}{HTML}{56645E}
\hypersetup{colorlinks=true,urlcolor=ink,linkcolor=ink,pdftitle={Prueba tecnica - Gestion de polizas},pdfauthor={Nicolas Ceron}}
\setlength{\parindent}{0pt}
\setlength{\parskip}{4pt}
\setlength{\emergencystretch}{2em}
\titleformat{\section}{\large\bfseries\color{ink}}{}{0pt}{}
\titlespacing*{\section}{0pt}{9pt}{4pt}
\setlist[enumerate]{leftmargin=1.5em,itemsep=2pt,topsep=3pt,parsep=0pt}
\pagestyle{fancy}\fancyhf{}
\fancyfoot[L]{\small\color{muted}Nicolas Ceron \;|\; Seguros Bolívar}
\fancyfoot[R]{\small\color{muted}\thepage\,/\,2}
\renewcommand{\headrulewidth}{0pt}
\newcommand{\tagline}[1]{\textbf{#1}\enspace}
\begin{document}
{\LARGE\bfseries\color{ink}Gestión de pólizas}\hfill{\small 16 septiembre 2026}\\[4pt]
{\large Prueba técnica \; / \; decisiones y entrega}\par
\href{https://github.com/nicoceron/seguros-bolivar-challenge/tree/assessment/final-2026-09-16}{\small GitHub: nicoceron/seguros-bolivar-challenge · rama assessment/final-2026-09-16}
\vspace{3pt}\hrule\vspace{3pt}

\section{1. Diseño del sistema}
\tagline{Tres patrones.} \textbf{Monolito modular}: pólizas y riesgos comparten transacción; evita coordinación distribuida innecesaria. \textbf{Hexagonal}: un puerto aísla el dominio del HTTP y de WebLogic. \textbf{Eventos con outbox}: cambio y envío pendiente se guardan juntos; una caída del CORE no pierde la intención de actualización.

\begin{center}\small\ttfamily
Cliente -> API Gateway -> Pólizas / Riesgos\\
\hspace*{5em}-> Base de datos + outbox\\
Outbox -> Adapter CORE -> WebLogic -> CORE legado\\
Outbox -> Broker -> Notificaciones -> Correo / SMS
\end{center}
{\small\color{muted}Implementado: API, base de datos, outbox y adapter al mock. Gateway, broker, notificaciones y WebLogic real pertenecen al diseño objetivo.}

\tagline{Modelo.} Póliza: tipo, estado, vigencia, meses iniciales, canon, prima, tomador y versión. Riesgo: póliza, inmueble, arrendatario y estado; relación 1:1 individual y 1:N colectiva. En producción, terceros con identidad y roles: individual, tomador = asegurado = arrendatario; colectiva, tomador = inmobiliaria/administración y beneficiario = arrendador por riesgo. La demo simplifica personas a nombres y un beneficiario por póliza.

\tagline{Operación objetivo.} Réplicas sin sesión; base de datos multi-zona, copias y restauración ensayada. Logs con correlación, sin claves ni datos personales; métricas de errores, latencia p95 y antigüedad del outbox. API pública \texttt{/v1}, contratos compatibles y migraciones graduales; la prueba conserva las rutas exigidas. Un planificador renovaría vencimientos con IPC registrado y unicidad póliza--vigencia, reutilizando el caso de uso. Creación y renovación publicarían eventos para notificar sin bloquear la API.

\section{2. Implementación y decisiones}
Java 21, Spring Boot, capas controller/service/repository, JPA y Flyway. Los seis endpoints exigidos y \texttt{/core-mock/evento} están implementados; crear y consultar por ID permiten probar sin depender del front. H2 facilita ejecución local; hay perfil PostgreSQL y Docker Compose.

\tagline{Reglas.} Individual: exactamente un riesgo al crear; agregar exige colectiva. Cancelar es baja lógica en cascada y repetirla no genera otro evento. Las mutaciones bloquean la póliza para evitar carreras entre agregar y cancelar. Renovar una cancelada devuelve 409; la nueva vigencia conserva los meses originales.
\[
 C'=\operatorname{redondear}_{2}\!\left(C(1+IPC/100)\right),\qquad P'=C'\,m.
\]
Se usa \texttt{BigDecimal}/\texttt{HALF\_UP}; derivar la prima del canon redondeado evita diferencias de centavos. IPC entra como porcentaje; sus límites y los demás supuestos están en README.

\tagline{CORE y seguridad.} Cada cambio confirma datos + outbox; el worker llama por HTTP al mock con \texttt{x-api-key: 123456}. Aplica timeouts, hasta cinco intentos y UUID estable. Un éxito de la API \emph{no} confirma el CORE. Los fallidos se conservan; producción requiere deduplicación, orden por póliza y conciliación. El mock solo registra el intento, no modifica un CORE real.

{\small\textbf{Verificación:} 55 pruebas aprobadas con \texttt{mvn verify}; ejecución: \texttt{mvn spring-boot:run}. Pruebas de reglas, API, rollback, concurrencia y HTTP del adapter. Renovación automática y notificaciones son diseño, no funcionalidades entregadas.}

\newpage
\section{3. Base de datos: tres estrategias}
\begin{enumerate}
\item \textbf{Índices:} evaluar \texttt{customers(country, customer\_id)} y \texttt{orders(customer\_id)}. Reducen búsquedas cuando México es selectivo; si abarca muchas filas, un \emph{hash join} con lectura completa puede ser preferible.
\item \textbf{Estadísticas y plan real:} actualizar con \texttt{DBMS\_STATS}; revisar sesgo de \texttt{country} e histogramas. Con estadísticas de ejecución habilitadas, usar \texttt{DBMS\_XPLAN.DISPLAY\_CURSOR} y \texttt{ALLSTATS LAST}; comparar filas estimadas/reales, lecturas y tiempo antes/después, sin imponer hints a ciegas.
\item \textbf{Materializar la consulta frecuente:} una vista materializada del resultado evita repetir la unión. Acordar frescura y coste de refresco; el refresco rápido depende de los requisitos de Oracle. Particionar por fecha no resuelve por sí solo esta consulta sin filtro temporal.
\end{enumerate}
No se suministraron datos ni plan de ejecución: estas son estrategias a medir, no mejoras cuantificadas.

\section{4. Git: traer únicamente el arreglo de seguridad}
Con el árbol de trabajo limpio, inspeccionar el commit que contiene el arreglo:
\begin{quote}\small\ttfamily
 git fetch origin\\
 git switch feature/new-login\\
 git show <SHA-del-arreglo>\\
 git cherry-pick -x <SHA-del-arreglo>
\end{quote}
\texttt{cherry-pick} aplica ese cambio, no el resto de \texttt{main}; \texttt{-x} conserva trazabilidad. Resolver conflictos, ejecutar pruebas, hacer \texttt{git add} de los archivos resueltos y \texttt{git cherry-pick --continue}; para desistir, \texttt{--abort}. Verificar dependencias del arreglo y elegir su commit, no un merge sin revisar sus padres.

\section{5. Liderazgo técnico}
\tagline{Cinco prioridades, primeras dos semanas.}
\begin{enumerate}
\item Días 1--2: clasificar los diez incidentes, contener los repetidos y asignar responsables con plan de reversión.
\item Días 1--3: acordar un alcance mínimo para la entrega de tres semanas, criterios de aceptación y riesgos visibles.
\item Semana 1: exigir revisión y CI; cubrir con pruebas los flujos que causaron incidentes antes de cambiarlos.
\item Semanas 1--2: priorizar deuda por impacto en fallos y entrega, no intentar eliminar todo el 40\% a la vez.
\item Desde el día 1: emparejar a cada junior con alguien experimentado, con tareas pequeñas y revisión diaria.
\end{enumerate}

\tagline{Organización de ocho personas.} Dos enfocadas en estabilidad y seis en la funcionalidad, incluidos los dos juniors acompañados. Rotar atención de incidentes, limitar trabajo en curso y mantener un único backlog priorizado. Revisar la distribución semanalmente según incidentes y avance.

\tagline{Métricas.} Línea base y revisión semanal: frecuencia de despliegue, tiempo desde cambio hasta producción, porcentaje de cambios fallidos y tiempo de recuperación. Complementar con incidentes repetidos, cumplimiento del objetivo de servicio y espera de revisión. No usar líneas de código como productividad individual.

\tagline{Prácticas obligatorias.} PR aprobado, CI verde, pruebas de reglas e integración, análisis de seguridad y secretos, migraciones compatibles, logs útiles y rollback probado. Toda corrección de incidente incluye una prueba de regresión.

\tagline{Presión del negocio.} Demostrar avances semanalmente y negociar alcance, no retirar validaciones. Entregar en incrementos detrás de banderas de funcionalidad; semana 3, despliegue gradual. Comunicar opciones y consecuencias; un riesgo crítico sin mitigación bloquea la salida.

\vfill
{\footnotesize\color{muted}\textbf{Base:} enunciado suministrado, módulos 1--5. \textbf{Referencias técnicas:}
\href{https://docs.spring.io/spring-data/jpa/reference/jpa/locking.html}{Spring Data JPA: locking};
\href{https://docs.oracle.com/en/database/oracle/oracle-database/19/arpls/DBMS_XPLAN.html}{Oracle: planes};
\href{https://docs.oracle.com/en/database/oracle/oracle-database/19/tgsql/optimizer-statistics-concepts.html}{estadísticas};
\href{https://docs.oracle.com/en/database/oracle/oracle-database/19/dwhsg/refreshing-materialized-views.html}{vistas materializadas};
\href{https://git-scm.com/docs/git-cherry-pick}{Git: cherry-pick}.}
\end{document}
